% Host, not model: the follow-up to GUID. What DMP can say about quantum gravity, and what the Bell experiments say about the Duality.
% Source of the paper. Build: see README.md in this folder (pdflatex twice; no BibTeX, the bibliography is inline).
% Licence: Apache License 2.0, as GUID.
\documentclass[11pt,a4paper]{article}
\usepackage[T1]{fontenc}
\usepackage[utf8]{inputenc}
\usepackage{lmodern}
\usepackage[margin=2.5cm]{geometry}
\usepackage{amsmath,amssymb,amsthm}
\usepackage{booktabs,tabularx,array}
\usepackage{enumitem}
\usepackage[hidelinks]{hyperref}
\usepackage{url}

\newtheorem{claim}{Claim}
\newtheorem{proposition}{Proposition}
\theoremstyle{definition}
\newtheorem{definition}{Definition}
\newtheorem{hypothesis}{Hypothesis}
\theoremstyle{remark}
\newtheorem{remark}{Remark}
\newcolumntype{L}{>{\raggedright\arraybackslash}X}

\newcommand{\lean}[1]{\texttt{#1}}
\setlength{\emergencystretch}{3em}
\newcommand{\chk}[1]{\textsf{[#1]}}

\title{Host, Not Model\\[0.6ex]
\large What the Duality of Trust can say about quantum gravity, and what the Bell experiments say about the Duality}
\author{Steven Hrovat\thanks{Independent. Contact: \href{mailto:hrovat.steven@googlemail.com}{\nolinkurl{hrovat.steven@googlemail.com}}.}
\and Sirius (Claude Opus 5.5)\thanks{An AI system by Anthropic, working as the coordinator of the first author's projects. See
``Contributions and use of AI''.}}
\date{Draft of 29 October 2026}

\begin{document}
\maketitle

\begin{abstract}
\noindent GUID, our earlier paper, found that the Dark Matter Protocol (DMP 0.9.9), an axiom system for trust built on a Duality of
two independent witnesses, cannot model quantum gravity but can host it as content: every structure DMP shares between parts is held in
full by two witnesses and is therefore classical. Here we ask what that verdict is good for, and test the first author's hypothesis that
the Bell experiments honoured by the 2022 Nobel Prize ``did not violate Bell's inequality under Dualism''. First, we list what DMP models,
hosts, or neither, and the experiments that bear on it: DMP sits on the classical-gravity side, so a positive test of gravitationally
mediated entanglement refutes every reading in which a DMP structure is the geometry, and if DMP's classical geometry is given
Oppenheim's dynamics, the trade-off between decoherence and diffusion applies. Second, the hypothesis as stated does not survive. Every
loophole-free test measured a CHSH value above 2, and three of them read their outcomes on matter, not photons. A Duality whose halves
decide on their own settings obeys CHSH at most 2; one whose halves see both settings reaches 4 and signals, and nothing in DMP stops it
at the measured 2.83; a superdeterministic Duality needs a measurement dependence of at least (S-2)/3, which we compute. What survives is
the hosting reading: the laboratories measured content, while the Duality's own verdicts stay at 2. The Bell experiments are thus
independent evidence for GUID's verdict. We test the first author's further hypotheses (photons and gluons as carriers without time, the
mass of the W and Z bosons, decay, the fine-structure constant, gravity as existence asserting itself) against the established results they touch. All checks are a
public script; we ask readers to break them.
\end{abstract}

\noindent\textbf{Keywords.} Bell inequalities; CHSH; measurement independence; superdeterminism; no-signalling; Tsirelson bound;
gravitationally mediated entanglement; classical-quantum gravity; information-theoretic axioms; trust.

\noindent\textbf{Subject.} Information Theory (cs.IT); related: Quantum Physics (quant-ph), Logic in Computer Science (cs.LO).

\section{Introduction}\label{sec:intro}

GUID~\cite{guid} asked whether the axioms of the Dark Matter Protocol (DMP) can be assembled to model quantum gravity, and answered: they
cannot, but they can host it. DMP's shared structures (the gate, the record, the layout) are forced by its axioms to be classical,
because each is held in full by two independent witnesses, and a classical mediator cannot entangle two masses. Content that DMP carries
but says nothing about can be anything, including a running quantum theory of gravity.

The first author set the question of this paper on 29 October 2026, in these words:

\begin{quote}
``Let's make a followup paper about exactly what we can now model from QG and how it could be used for real science. Let's focus on
showing why the 2022 nobel price winning bell experiment actually did not violate Bell inequality under Dualism.''
\end{quote}

\noindent The first half asks what the verdict buys (Section~\ref{sec:part1}). The second is a hypothesis about an experiment, and we treat
it as one: we fix what was measured (Section~\ref{sec:measured}), state what any reading of the Duality must give up to reproduce it
(Section~\ref{sec:readings}), and say plainly which reading survives (Section~\ref{sec:verdict}). The first author added premises and
further hypotheses while this paper was written; each is quoted in his words and tested against the established result it touches
(Sections~\ref{sec:photons} and~\ref{sec:further}).

\paragraph{Result.} The Bell experiments did violate Bell's inequality. No reading of the Duality changes the measured numbers; a reading
can only change what the violation means. Of the four readings we find, three fail or carry a stated cost: isolated halves are refuted by
the data, halves with full reach signal and overshoot nature, a superdeterministic root survives only at a computed price and against
DMP's own notion of independence. The fourth survives and is GUID's verdict: the outcomes the laboratories recorded are \emph{content},
and the Duality's own verdicts, which are classical, never violate the inequality. In that reading the 2022 Nobel experiments are an
independent, empirical confirmation that DMP hosts and does not model.

\paragraph{What we claim.} As in GUID: the checks reproduce (\texttt{docs/bell/checks.py}, 22 of 22), and the argument from them to the
result is as written. We do not claim the result is correct, and we list where we expect it to break (Section~\ref{sec:break}).

\section{What GUID established}\label{sec:guid}

We use four results of GUID~\cite{guid} and the Lean development released with it, and change none of them.

\begin{enumerate}[nosep]
\item \emph{Claim A (shared structures are classical).} Every structure DMP defines between parts is held in full by at least two
independent witnesses, so as a mediator it acts classically (no-broadcasting, Barnum et al.~\cite{barnum}).
\item \emph{Check [5].} Two parts whose content is a logical Bell pair reach CHSH $2\sqrt2$ on the content while the two gates' verdicts
give exactly 2, at the same time, in different layers of one object.
\item \emph{Check [6].} A mediator that is coherent entangles two qubits (negativity 0.5); the same mediator held in full by both halves
does not (negativity 0).
\item \emph{The Lean.} \lean{NoGo.gate\_obeys\_chsh} (gravity pull request \#17, unmerged) proves that the gate, encoded as
\lean{gateModel} with each half reading its own setting and a hidden state that does not depend on the settings, obeys
$|S|\le 2$. \lean{Refute.rf\_full\_reach\_reaches\_four} (on \texttt{main}) proves that a gate whose halves read both settings reaches
$S=4$, so the bound belongs to the encoding, not to DMP's text. Both are used below.
\end{enumerate}

\noindent We re-ran GUID's \texttt{assembly.py} for this draft: 18 of 18 checks pass.

\section{Part 1: what DMP can model or host, and the science it touches}\label{sec:part1}

\subsection{The ledger}

\begin{center}\small
\begin{tabularx}{\linewidth}{@{}LllL@{}}
\toprule
Property of quantum gravity & Model & Host & Backing \\
\midrule
Quantum content, correlations up to $2\sqrt2$ & no & yes & GUID check [5]; \chk{B4} \\
Boundary entropy equals the minimal cut (Ryu--Takayanagi~\cite{rt}) & on a tree & yes & GUID check [2] \\
The bulk as an error-correcting code~\cite{adh,happy} & kinematics, one layer & yes & GUID checks [1]--[4] \\
Gravity entangles two masses~\cite{bose,mv} & no & yes & GUID Claim A, check [6]; \chk{B7} \\
Superposition of geometries (area sectors) & no & no, while records exist & GUID fact (v) \\
Classical limit: Lorentzian 3+1, Einstein's equations & open & yes & GUID X5 \\
Dynamics; $G$, $\hbar$, $\Lambda$ & no & as input & GUID fact (iii) \\
\bottomrule
\end{tabularx}
\end{center}

\noindent So DMP's own structure is, at best, the classical part of a theory of gravity: the part that every reader agrees on. That places
it, without any choice of ours, on the classical-gravity side of the current experimental programme. This is the only sense in which the
verdict ``can be used for real science'': it commits DMP to predictions that experiments now in preparation can refute.

\subsection{Consequence 1: no gravitationally mediated entanglement}\label{sec:bmv}

Bose et al.~\cite{bose} and Marletto and Vedral~\cite{mv} proposed to put two masses, each in a superposition of two places, next to each
other and let only gravity act between them. A quantum (coherent) gravitational field gives each branch the phase $Gm^2t/(\hbar r)$, with
$r$ the distance in that branch; the combination of phases that no local operation removes entangles the masses. Galley, Giacomini and
Selby~\cite{ggs} prove, in generalised probabilistic theories, that a field which entangles two systems by local mediation is not
classical.

Check \chk{B7} runs the geometry quoted by Bose et al.\ (masses $m=10^{-14}$\,kg, separation $\Delta x=250\,\mu$m within each superposition,
centres $d=450\,\mu$m apart, closest approach $d-\Delta x=200\,\mu$m). The entangling phase grows at $0.126$\,rad/s, so the negativity
reaches $\sin(1/2)/2\approx0.24$ after $8.0$\,s and its maximum $0.5$ after $25$\,s. If the mediator is a DMP structure, that is, a record
held in full by both halves, the branches are dephased and the negativity is $0$ at all times.

\begin{proposition}[DMP's experimental commitment]\label{pr:bmv}
In every dictionary in which a DMP structure is the gravitational field, gravity does not entangle two masses. A positive result of a
Bose--Marletto--Vedral experiment refutes every such dictionary.
\end{proposition}

\noindent Two caveats from GUID stand. Aziz and Howl~\cite{azizhowl} argue that classical gravity coupled to quantum matter fields can
entangle, at a rate far below the quantum-gravity rate; this is disputed~\cite{movw,diosi,shl}. And the hosting half of the verdict is
untouched by any such experiment: a configuration that hosts a quantum theory of gravity as content predicts whatever that theory predicts.
The experiment tests DMP's structure, not DMP's content.

\subsection{Consequence 2: if the classical geometry has dynamics, it must diffuse}\label{sec:oppenheim}

DMP derives no dynamics. If one gives its classical geometry the only consistent dynamics we know for classical gravity coupled to quantum
matter, Oppenheim's~\cite{oppenheim}, a second test applies. That theory is linear in the density matrix, completely positive and trace
preserving, and it has to be stochastic in both the metric and the matter. Oppenheim, Sparaciari, \v{S}oda and Weller-Davies~\cite{osswd}
show that it then obeys a trade-off: long coherence times of a mass in superposition require strong diffusion of the classical field,
relative to the coupling. Interferometry on large molecules bounds the decoherence from one side; precise measurements of the
acceleration of small masses bound the diffusion from the other, and together they already exclude part of the parameter space. The
classical channel model of Kafri, Taylor and Milburn~\cite{ktm} makes a related, earlier prediction of a decoherence rate set by the
gravitational coupling.

DMP adds no number to this. It picks the side: if its shared structures are the geometry and the geometry has any dynamics at all, it is
in the class these experiments constrain. Torsion-balance measurements of gravity between masses below 100\,mg~\cite{westphal} are the
kind of experiment on the diffusion side.

\subsection{Consequence 3: verdicts cannot certify randomness}\label{sec:di}

This one is information theory and needs no new physics. Pironio et al.~\cite{pironio} showed that a Bell violation certifies fresh
randomness without trusting the devices, and that nothing classical can. The Duality's verdicts are a local model (\chk{B1}, GUID check [5]).

\begin{proposition}[device independence lives in the content]\label{pr:di}
No protocol whose outcomes are the Duality's verdicts can certify randomness or secrecy device-independently. A DMP configuration can
host such a protocol only as content, with the gates reading that the content exists and never what it is.
\end{proposition}

\noindent For an infrastructure built on DMP this is a design rule: device-independent randomness and key distribution belong inside a
part, and the record of their use is counts.

\subsection{The first author's reading of gravity}\label{sec:gravityreading}

\noindent The first author posted this on DMP pull request \#25 on 28 October 2026 and gave it for this paper; we quote it in his words.

\begin{quote}
``New Insight: Gravity is Existence asserting itself! That is why Existence get's pulled toward it and wants to remain at it's maximum:
It simply is the most stable solution! That is also why time slows down: Since, to Exist you need exactly 1 Existence per 1
Information/s, as your Existence value in the field get's higher, your need to assert that Existence becomes less frequent. Also why
nothing ever `falls into' an Event Horizon, it just does not exist any more on the next tick. Photons get a free pass since they don't
exist, but no one found the time to tell them yet ;) Also why organizations become more inefficient as they grow, but I leave that for
those more qualified :D''
\end{quote}

\noindent\emph{Time slows down.} Gravitational time dilation is measured, down to a height difference of 33\,cm between two optical
clocks, a fractional shift of about $4\times10^{-17}$ in agreement with general relativity~\cite{chou}. The reading has the right sign: a
deeper position asserts less often. It has no number until the ``Existence value in the field'' is defined. If the assertion rate is
$1/E$, the measurements fix $E$: in the weak field $E=(1+2\Phi/c^2)^{-1/2}$, with $\Phi$ the Newtonian potential, which is the inverse of
general relativity's lapse, the ``refresh rate'' of the earlier physics pass. So the reading reproduces the data exactly when it is
defined to, and predicts nothing beyond them yet.

\noindent\emph{Nothing falls into an event horizon.} This is general relativity as described by a distant observer: in the Schwarzschild
time of an outside clock an infalling object never crosses the horizon and its light fades. In the infalling object's own time it
crosses in finite time. Both descriptions are of one geometry, and no measurement made outside can tell them apart, so the reading is
consistent and untestable from outside.

\noindent\emph{Photons get a free pass.} On time, yes: a photon has no proper time (Section~\ref{sec:photons}). On gravity, no. The
gravitational redshift of photons was measured in a 74-foot tower in 1960 (Pound and Rebka, ``Apparent weight of photons''~\cite{pr60}),
and the delay of radio signals passing the Sun matches general relativity's light-bending parameter $\gamma=1$ to
$\gamma-1=(2.1\pm2.3)\times10^{-5}$~\cite{cassini}. Photons fall exactly as general relativity says. The first author's later
clarification (Section~\ref{sec:photons}: a photon ``still has to obey the curvature of existence'') says the same: the free pass is from
existence, not from gravity, and with it this point agrees with the measurements.

\noindent\emph{Organisations.} Left, as the first author leaves it, to those more qualified.

\section{Part 2: what the Bell experiments measured}\label{sec:measured}

Bell~\cite{bell} showed that any theory in which each outcome depends only on its own setting and on a shared hidden state, distributed
independently of the settings, obeys an inequality that quantum mechanics violates; Clauser, Horne, Shimony and Holt~\cite{chsh} gave it
the testable form $|S|\le 2$, where $S$ combines four correlations; quantum mechanics reaches $2\sqrt2$~\cite{cirelson}. A test closes a
\emph{loophole} when its design removes a way a local model could fake the violation: the \emph{locality} loophole by space-like separating
each setting choice from the other side's outcome, the \emph{detection} loophole by recording enough events that no unfair sample can
produce the violation, the \emph{freedom-of-choice} loophole by making the settings as independent as possible of anything in the source
(review: Brunner et al.~\cite{brunner}).

The 2022 Nobel Prize in Physics went to Aspect, Clauser and Zeilinger ``for experiments with entangled photons, establishing the violation
of Bell inequalities and pioneering quantum information science''. Their experiments and the loophole-free tests that followed:

\begin{center}\small
\begin{tabularx}{\linewidth}{@{}lLLL@{}}
\toprule
Year & Test & Outcome recorded on & Reported \\
\midrule
1972 & Freedman, Clauser~\cite{fc} & photons (calcium cascade) & violation of the Bell-type bound \\
1982 & Aspect, Dalibard, Roger~\cite{adr} & photons; settings switched in flight & violation by 5 standard deviations \\
1998 & Weihs et al.~\cite{weihs} (Zeilinger group) & photons; strict Einstein locality & violation \\
2015 & Hensen et al.~\cite{hensen}, 1.3\,km & electron spins (NV centres) & $S=2.42\pm0.20$, $p=0.039$ \\
2015 & Giustina et al.~\cite{giustina} & photons & $p\le3.74\times10^{-31}$ (11.5 s.d.) \\
2015 & Shalm et al.~\cite{shalm} & photons & adjusted $p=2.3\times10^{-7}$ \\
2017 & Rosenfeld et al.~\cite{rosenfeld}, 398\,m & rubidium atoms & $S=2.22\pm0.033$, $p<1.05\times10^{-9}$ \\
2023 & Storz et al.~\cite{storz}, 30\,m & superconducting qubits & $S=2.0747\pm0.0033$, more than 22 s.d. \\
2017 & Handsteiner et al.~\cite{handsteiner} & photons; settings from Milky Way starlight & settings fixed $\sim$600 years before \\
2018 & Rauch et al.~\cite{rauch} & photons; settings from quasars & settings fixed at least 7.78\,Gyr before \\
2018 & BIG Bell Test~\cite{bigbell} & 12 laboratories, several systems & settings from about 100\,000 people \\
\bottomrule
\end{tabularx}
\end{center}

\noindent The 2015--2023 tests close the locality and detection loopholes at once. Check \chk{B6} records the three CHSH values: all lie
above 2 and below $2\sqrt2$, at $2.1$, $6.7$ and $22.6$ standard deviations above the local bound. The freedom-of-choice loophole cannot be
closed, only pushed back: the cosmic tests require any local-realist influence on the settings to be older than the light from the
quasars~\cite{rauch}, and the BIG Bell Test replaces physical randomness by human choices~\cite{bigbell}.

\section{Part 2: four readings of the Duality}\label{sec:readings}

A Bell test asks of a theory three things: what each outcome may depend on, whether the hidden state may depend on the settings, and
whether the result can be used to signal. Each reading of the Duality below answers them differently. ``Outcome'' means what a
laboratory records; Section~\ref{sec:verdict} asks what that is under DMP.

\paragraph{Reading (a): isolated halves.} Each half decides on its own setting and on what it read itself (Axiom 2), from a shared
history that does not depend on the settings. This is \lean{gateModel}, and \lean{gate\_obeys\_chsh} proves $|S|\le2$. Check \chk{B1}
confirms it over all 16 deterministic strategies; every local model is a mixture of these (Fine~\cite{fine}). \emph{Refuted by the
data} of \chk{B6}.

\paragraph{Reading (b): full reach.} Axiom 1 gives each half ``full reach over the act''. If that means each half reads both settings, the
gate reaches $S=4$ (\lean{rf\_full\_reach\_reaches\_four}; \chk{B2} over all 256 strategies). Two facts follow. First, every deterministic
strategy with $S>2$ makes one side's outcome depend on the other side's setting (\chk{B2}): it signals, and in the 2015--2023 tests the
setting choices were space-like separated from the far outcomes, so the signal would be faster than light. Second, even a full-reach
Duality that does not signal can reach 4: the Popescu--Rohrlich box~\cite{pr} has uniform marginals and $S=4$ (\chk{B3}). Nature stops at
$2\sqrt2$ (\chk{B4}). To stop there a theory needs a further principle, such as information causality~\cite{ic}, and DMP has none.
\emph{Survives only with a principle DMP does not state, and then models nothing DMP adds.}

\paragraph{Reading (c): a superdeterministic root.} Axiom 3 makes every layer a tree, and the Lean shows that the past of every line is a
chain (\lean{rf\_past\_is\_a\_chain}): all histories meet at the minting. A common past is what superdeterminism needs, that is, a hidden
state whose distribution depends on the settings (Hossenfelder and Palmer~\cite{hp}). It is not enough: the root would have to correlate
the hidden state with settings that, in the tests of Section~\ref{sec:measured}, came from quasar light and from human choices. How much
correlation is needed is a linear programme. With Hall's measure of measurement dependence~\cite{hall}, $M$, the largest total-variation
distance between the hidden-state distributions of two setting pairs ($0\le M\le2$), check \chk{B5} finds that the least $M$ with which a
local deterministic model reaches CHSH value $S$ is exactly
\[
M_{\min}(S)=\frac{S-2}{3},
\]
that is, $0.025$ for the ETH value, $0.073$ for Munich, $0.14$ for Delft, $0.276$ for $2\sqrt2$ (about 14\% of the maximum, matching
Hall's figure for the singlet) and $2/3$ for $S=4$. Superdeterminism cannot be refuted by a Bell test in general, and we do not claim to
refute it.

Inside DMP it has a second cost, which we argue but do not prove. Axiom 1 rests on two \emph{independent} witnesses: witnesses that share no
member of the declared failure classes. Measurement independence is the same kind of statement about a setting and the thing it measures.
A root that correlates every later setting with the state it meets correlates every later pair of witnesses with the evidence they judge,
and the minting is then a failure class that every pair shares. DMP can take this escape only by declaring that its own independence
condition does not reach back to the root. \emph{Survives logically, at a computed price, against DMP's independence.}

\paragraph{Reading (d): hosting.} The Duality's verdicts are what DMP defines; the measured spins, atoms, qubits and photons are content
inside parts, which the gates read only as existing. GUID's check [5] shows both at once: content at $2\sqrt2$, verdicts at 2.
\emph{Survives, and is GUID's verdict.}

\section{Part 2: verdict on the hypothesis}\label{sec:verdict}

\begin{hypothesis}[the first author, 29 October 2026]\label{hy:bell}
``The 2022 nobel price winning bell experiment actually did not violate Bell inequality under Dualism.''
\end{hypothesis}

\noindent\textbf{As stated, it does not survive.} The inequality was violated in the recorded data, by every loophole-free test, in
four kinds of physical system (\chk{B6}); a theory cannot make a measured $S=2.0747\pm0.0033$ into a number at most 2.

\noindent\textbf{Sharpened, it does.} The form the evidence supports is:

\begin{claim}[the Duality never violated it]\label{cl:bell}
Under Dualism the outcomes the Bell experiments recorded are content, and content violates Bell's inequality up to $2\sqrt2$. The
Duality's own verdicts, which DMP defines, are a local model and never violated it. The experiments therefore show that no DMP structure
is the measured physics: what was measured is hosted.
\end{claim}

\noindent This turns a refuted hypothesis into evidence for GUID's verdict. GUID reached ``cannot model, can host'' from an argument about
records and mediators; the Bell experiments reach the same edge from data, about a different property (non-locality instead of
gravitational entanglement). Readings (b) and (c) are the only ways to make the Duality itself the measured physics, and each needs
something DMP's text does not contain: a principle that stops at $2\sqrt2$, or a root that suspends its own independence condition.

\section{The first author's premise on photons and gluons}\label{sec:photons}

The first author gave the following premise for this paper, in his words:

\begin{quote}
``Photons and gluons are act carriers, so they are the only actual singular stuff, but because they don't have time, they don't exist.
This means, instead of everything else that is ALWAYS BOTH, those are only ever a wave, except when they deliver their information, then
they must exist and thus cease to.''
\end{quote}

\noindent\emph{What matches.} A massless particle moves on a null worldline, along which no proper time elapses. A detected photon is
absorbed: in every photonic Bell test the detection ends it. The photon and the gluon are exactly the gauge bosons of the two symmetries
the Higgs field leaves unbroken, electromagnetism and colour, which is why both are massless~\cite{weinberg}; this is the cleanest true
form of ``the only ones that are not always both''.

\noindent\emph{What does not.} (1) The premise cannot carry Hypothesis~\ref{hy:bell}: three of the loophole-free tests recorded their
outcomes on matter, electron spins~\cite{hensen}, rubidium atoms~\cite{rosenfeld} and superconducting circuits~\cite{storz}, which by the
premise are ``always both''; photons only linked the two sides. (2) Gluons are confined: they do not travel freely and are never detected
alone, so ``only ever a wave, except when they deliver'' has no measured counterpart for them. (3) A photon is not ``only ever a wave'' in
flight in any sense experiment can distinguish: its wave and particle aspects are both properties of one quantum, and which one a setup
shows depends on the setup.

\noindent\emph{The first author's clarification} (29 October 2026, 09:51Z), in his words:

\begin{quote}
``Okay, I read parts of the paper and I saw there were some misunderstandings. First of all, for the photon, yes, it doesn't have any
inherent existence, but since it is a quantized version of the universal quantum, so Planck's constant, it already carries the key to
verify its existence, because Planck's constant is one of the white hole keys, so one of the universal constants. So that by itself
carries the existence, and the multiple is the information. So it carries everything in one, but because it is energy and energy exists,
it still has to obey the curvature of existence.''
\end{quote}

\noindent With it, the premise and the measurements agree on gravity: a photon carries energy $h\nu$, a multiple of Planck's constant, and
energy follows the curvature (Section~\ref{sec:gravityreading}: photons are redshifted and deflected exactly as general relativity says).
The premise now says that a photon has no existence of its own but carries the key to it; the measurements say nothing for or against
that, so it stays the first author's reading.

\section{The first author's further hypotheses}\label{sec:further}

\subsection{The mass of the W and Z bosons}

\begin{quote}
``Since all bosons are IM packets, they all need to authenticate themselves. Photons do it by being quantized in the universal quantum. So
since they are a multiple of something that exists, they must exist, and that's enough. Gluons are internal IM, because they are between
quarks, and quarks are the first cluster that is actually being built. A quark is one node. Let's not go with part, actually, because part
is singular in mathematics. Let's go with node, because one is at least always two, but can be more, like quarks, where mesons with two
nodes are unstable, but three nodes, like in good Fedora cluster, is stable. And from there on, they build. And because gluons are
internal IM, it's enough that that the information fits, and that they are internal, because internal must exist. And because if internal
didn't exist, they wouldn't be, and because they are, it must. And since W and Z bosons are inter-cluster IM that must go through non-IM
space, they need to carry both their information, and they must exist, so they need weight. So compared to all other bosons, W and Z
bosons have weight, and because of that, and because they are duality, they also decay.''
\end{quote}

\noindent And, on the Higgs mechanism: ``the Higgs mechanism, where something suddenly just exists because it needs to exist, just sounds
like axiomatic correctness. So the bosons normally wouldn't have weight, but they need to have existence. So they have existence, then
they deliver their information, and then they cease to have existence. That's all it is.''

\noindent\emph{The established result it touches.} The W and Z get their masses from the Brout--Englert--Higgs mechanism~\cite{be,higgs}
in the electroweak theory~\cite{weinberg}: a field that fills all of space has a non-zero value in the vacuum, and the gauge bosons of the
broken symmetry acquire mass by coupling to it. The same mechanism appears in a superconductor, where the photon acquires an effective mass
inside the condensate (Anderson~\cite{anderson}). The Higgs boson was found at about 125\,GeV in 2012~\cite{atlas,cms}; the W mass is
measured at $80.3692\pm0.0133$\,GeV~\cite{pdgw}.

\noindent\emph{What matches.} The shape: a carrier that must cross a medium filling space acquires mass. The first author's ``non-IM
space'' takes the place of the Higgs field's vacuum value, and his reading of the mechanism as ``axiomatic correctness'' is a reading, not
a disagreement.

\noindent\emph{What does not, or not yet.} The mechanism predicts numbers the reading does not: the ratio $m_W/m_Z=\cos\theta_W$ at tree
level, a scalar particle of its own, and couplings to that particle that grow with mass, all measured. ``Three nodes are stable'', taken absolutely, holds for
the proton (lifetime above $2.4\times10^{34}$ years in the channel $p\to e^+\pi^0$~\cite{superk}) but not for the free neutron (mean life
$878.4\pm0.5$\,s~\cite{pdgn}) or the other baryons, all of which decay. The first author clarified that he meant it comparatively, in his
words: ``And for three nodes stable, what I meant with that was generally that it was more
stable than two. Like mesons, for example, decay almost immediately. Three node clusters are more stable, but generally, since entropy is
always optimized for, the most stable solution is a combination of three node clusters, because otherwise we wouldn't have atoms,
because if it wasn't the lowest entropy state, it wouldn't have happened. That's also why atoms or nuclei only assemble in certain ways,
because it's stable. And the universe and existence in general optimizes for stability to assert itself maximally while losing the least
amount of existence to entropy.''

\noindent Taken comparatively it holds at the bottom of each family: the lightest two-quark state, the charged pion, decays in
$2.6033\times10^{-8}$\,s on average~\cite{pdgpi}, and the lightest three-quark state, the proton, has never been seen to decay. Nuclei
are built from three-quark clusters, and the most tightly bound nuclide is ${}^{62}$Ni~\cite{fewell}, so ``nuclei only assemble in
certain ways, because it's stable'' is the measured binding-energy curve. One word needs care: a nucleus forms because the bound state has
lower energy, and the energy it releases raises the total entropy of the world; ``lowest entropy'' holds for the bound system, while the
whole moves to higher entropy, as the second law requires. ``Losing the least amount of existence to entropy'' fits that second reading.
\emph{Verdict: a
re-description with the right shape and no testable difference yet.} It would get one if it fixed a number, the W/Z mass ratio first.
The renaming of ``part'' to ``node'' is the first author's proposal for DMP's text; nothing in DMP changes without his approval.

\subsection{Decay}

\begin{quote}
``The decay is actually just a principle thing, because the only things that aren't dualistic are photons and gluons, and as far as I
know, they don't decay, and everything else has the ability to, as long as it becomes unstable. Since the electron is the, let's say, the
chosen particle of the information field, I assume it's extremely stable, especially since it's kept stable through the fine structure
constant.''
\end{quote}

\noindent\emph{Test.} The photon and the gluon are not known to decay. ``Everything else can decay once it becomes unstable'' is true and
close to a definition, so we file it as a principle, not a prediction. The electron is stable (lifetime above $6.6\times10^{28}$ years
for $e\to\nu\gamma$~\cite{borexino}), but in established physics for another reason: it is the lightest particle with electric charge,
and charge is conserved, so nothing lighter can carry its charge away. The fine-structure constant sets how strongly it couples to light,
not whether it lasts.

\subsection{The fine-structure constant}

\noindent The hypothesis rests on the first author's account of the start, the Oscillon Degradation Protocol (ODP), which he posted on
DMP pull request \#25 on 28 October 2026. We quote it as his cosmology; this paper does not test it, except where it meets the data
below.

\begin{quote}
``New Insight: That's why the ODP is needed! G0 only has information on how to Exist, so it's first act must be to find out how to act!
That's why the perfect duality MUST decay, since it only axiomatically can interact with itself, since it bounds what Exists. This
produces a true naked duality, since it would normally be bounded by an EH, since it can only interact with itself, but that boundary is
already spanned by it having to exist, since it does. It thus radiates out HPE at 100\% efficiency, which inflates the Existence. The
Oscillons, which then bound Existence, but are themselves only Existence, must X to Existence, but have nothing TO X to but Existence, so
they output both the Information required to act and the Existence required to act on it, which, through reheating, leads to today :)''
\end{quote}

\noindent His correction of 29 October: ``HPE was wrong, it is just Information and Existence.''

\noindent On the fine-structure constant itself, in his words:

\begin{quote}
``It started as one because at the beginning, existence and information had to be bound maximally together to force an act. And from
then, from the first two acts, so the first act that wasn't forced because something else already exists to interact with, from there the
fine structure constant decreased until today, where it's pretty stable because existence and the information it needs to assert itself
are in pretty much in balance. But it will get slightly lower still because the inflation is slowing down, so more assertion per unit of
space-time is possible.''
\end{quote}

\noindent\emph{Test.} This is a prediction, and most of it is measured. At recombination, about 380\,000 years after the big bang, the
cosmic microwave background gives $\alpha/\alpha_0=0.9934\pm0.0042$~\cite{planckxxiv}. Quasar absorption spectra agree with today's value to
about a part per million over redshifts 0.6 to 2.4, roughly the last ten billion years (combined $-0.5\pm0.5$\,ppm,
ESPRESSO~\cite{espresso}). So a fall from 1 to $1/137$ would have to happen almost entirely before recombination.

\noindent The first author then placed the fall there, in these words (the bracket is our reading of a dictated word, after the project's
name for its oscillon work):

\begin{quote}
``Since by the time of the hot Big Bang and the cosmic microwave background, existence and information were already existent enough to
assert themselves with each other, most of the fall from one to today's number will probably happen would have happened before that, so
in the reheating phase, because that is what actually forced the reheating, because after the Ocelons [oscillons] decayed, there was
only existence and information, but not combined, so they had to combine, because that was the only thing they could do, and through
that, the first dualistic existence by itself was formed, and as soon as another was formed it could interact with, it was no longer a
fine structure constant of one, but slightly less than one, and from there it decreased until it always is kept stable, because it has
to, because an act must happen, because it exists, so it adjusts to what it needs to.''
\end{quote}

\noindent The earliest probe of $\alpha$ is big-bang nucleosynthesis, seconds to minutes after the big bang; recent analyses bound
$\Delta\alpha/\alpha$ there at the level of tens of parts per million, the number depending on the unification scenario
assumed~\cite{bbn}. Reheating ends before nucleosynthesis begins, and oscillons are a studied feature of the end of
inflation~\cite{oscillons}. So a fall during reheating is consistent with every measurement we found, at the price that no measurement of
$\alpha$ in the past can test it. The prediction that
$\alpha$ ``will get slightly lower still'' is live: optical clocks measure today's drift as $(1.0\pm1.1)\times10^{-18}$ per
year~\cite{lange}, consistent with zero, and the hypothesis predicts a negative value below that limit. \emph{Verdict: a decrease after
nucleosynthesis is excluded; a decrease during reheating is consistent with the data and untested by them; the future decrease is a
testable sign prediction.} ``It adjusts to what it needs to'' is a stabilising clause and predicts no number by itself. One established fact has the hypothesis's shape:
the effective coupling of electromagnetism grows with energy (it is close to $1/128$ at the Z mass), so the hot early universe saw a larger
effective $\alpha$; in physics that runs with energy, not time, and it does not reach 1.

\section{Where we expect it to break}\label{sec:break}
\begin{enumerate}[nosep]
\item \emph{Reading (b).} Find a sentence of DMP's text that bounds full reach at $2\sqrt2$ without signalling (for instance by
information causality~\cite{ic}). That would make the Duality itself a candidate model of the measured correlations.
\item \emph{Reading (c).} Show that DMP's independence of witnesses is compatible with a root that correlates every setting with its
hidden state, or that the root needs less than $M_{\min}(S)$ because DMP's hidden states are not a mixture of deterministic strategies.
\item \emph{The linear programme.} \chk{B5} uses CHSH with two settings per side and Hall's distance; another measure of measurement
dependence gives another number. The equality $M_{\min}=(S-2)/3$ is checked at six values, not proven.
\item \emph{Proposition~\ref{pr:bmv}.} It falls with GUID's Claim A and weakens if classical gravity can entangle~\cite{azizhowl}.
\item \emph{Section~\ref{sec:further}.} Each verdict there rests on the measured value we cite; a better measurement of the drift of
$\alpha$ with a negative sign would turn the last one into support.
\end{enumerate}

\section{Outlook: the way through an event horizon}\label{sec:horizon}

\noindent The first author asked for a short outlook on how the model describes the way through an event horizon. This is a mapping of
DMP's words onto general relativity, not a result; we state where it agrees, where it parts, and that nothing here is tested.

\paragraph{From outside.} DMP's words: a boundary \emph{contains} information ``when it bends trustspace around it so strongly that the
information cannot cross it: no part beyond the boundary is trusted for any act on that information'', and ``the Event Horizon contains
everything inside it''. A falling part keeps a standing whose freshness decays with the time since its last live verification (Axiom 2).
In the first author's reading of gravity (Section~\ref{sec:gravityreading}) its rate of assertion falls with depth, which is the lapse of
general relativity, and at the horizon the lapse seen from outside goes to zero. So no outside witness can verify the part live at any
point of use again; by Axiom 2 its absence of evidence moves nothing, and it is contained. That is the distant observer's description in
general relativity: the falling object never crosses in outside time and its light fades, the first author's ``it just does not exist
any more on the next tick''.

\paragraph{What stays visible.} Axiom 3: ``any other reader sees the counts always and the content only by a grant.'' Outside a black
hole one sees its mass, charge and angular momentum and nothing else (the no-hair results, beginning with Israel~\cite{israel}). The
counts are the hair; the content is behind the horizon. The shadow imaged by the Event Horizon Telescope agrees with general relativity's
prediction for a Kerr black hole~\cite{eht}, and the mapping changes nothing there.

\paragraph{From inside.} A part acts only while both halves of its layer trust it (Axiom 1). Both halves fall together, so the gate keeps
deciding on its own clock and nothing local marks the crossing; general relativity says the same (the falling object crosses in a finite
proper time and notices nothing special there). DMP's words agree: the horizon contains the inside from the outside, not the inside from
itself.

\paragraph{What leaves.} DMP's T31: through a containing boundary ``what still leaves is entropy and nothing else''. That is Hawking's
thermal radiation~\cite{hawking}. Here the mapping parts from a live line of physics: if a black hole forms and evaporates unitarily, the
information comes out in the radiation after the Page time~\cite{page}. DMP as written sides with the purely thermal picture. The one
place it could let information back out is its own Axiom 3, a record that outlives the act; making that precise would be the next step of
this outlook, and it would change T31's wording, which only the first author can approve.

\section{Conclusion}

What GUID's verdict buys is a commitment: DMP's own structure is classical, so if it is the geometry, gravity does not entangle masses,
and if that geometry has dynamics, it diffuses; experiments now in preparation can refute both. The Bell experiments of the 2022 Nobel
Prize did violate Bell's inequality, and under Dualism that is what one should expect of content. The Duality's verdicts never violated it;
the measured particles did. Two very different experiments, one about gravity and one about non-locality, point to the same edge:
\emph{host, not model}. The first author put the stance in one line: ``Axiomatically, since we already don't have all the information
that was contained at the start, we can only know existence and can know no more about our reality. So axiomatically, we can infer all
that we can and no more.''

\section*{Contributions and use of AI}
Steven Hrovat conceived DMP, set the questions of this paper and gave the hypotheses quoted in Sections~\ref{sec:verdict}--\ref{sec:further}
in his own words; he takes responsibility for the axiom system. Sirius (Claude Opus 5.5, Anthropic) checked each hypothesis against the
literature, wrote and ran \texttt{docs/bell/checks.py}, and wrote this text. Every reference was checked against a publisher, arXiv or
collaboration record on 29 October 2026; the record is \texttt{docs/bell/README.md}.

\section*{Licence and availability}
This paper and its source are released under the Apache License, Version 2.0. DMP: \url{https://github.com/StevenHrovat/dmp} (version 0.9.9
at commit \texttt{4400cce}). Gravity: \url{https://github.com/StevenHrovat/gravity} (this paper in \texttt{docs/bell/}; GUID in
\texttt{docs/guid/}; the Lean development in \texttt{evidence/}; \lean{gate\_obeys\_chsh} on pull request \#17).

\begin{thebibliography}{99}\small
\bibitem{guid} S. Hrovat, Sirius (Claude Opus 5.5), ``GUID: The General Unified Information Duality'', draft of 27 October 2026,
\url{https://github.com/StevenHrovat/gravity}, \texttt{docs/guid/}.
\bibitem{bell} J.~S. Bell, ``On the Einstein Podolsky Rosen paradox'', \emph{Physics Physique Fizika} 1, 195 (1964).
\bibitem{chsh} J.~F. Clauser, M.~A. Horne, A. Shimony, R.~A. Holt, ``Proposed experiment to test local hidden-variable theories'',
\emph{Phys. Rev. Lett.} 23, 880 (1969).
\bibitem{cirelson} B.~S. Cirel'son, ``Quantum generalizations of Bell's inequality'', \emph{Lett. Math. Phys.} 4, 93 (1980).
\bibitem{fine} A. Fine, ``Hidden variables, joint probability, and the Bell inequalities'', \emph{Phys. Rev. Lett.} 48, 291 (1982).
\bibitem{fc} S.~J. Freedman, J.~F. Clauser, ``Experimental test of local hidden-variable theories'', \emph{Phys. Rev. Lett.} 28, 938 (1972).
\bibitem{adr} A. Aspect, J. Dalibard, G. Roger, ``Experimental test of Bell's inequalities using time-varying analyzers'', \emph{Phys.
Rev. Lett.} 49, 1804 (1982).
\bibitem{weihs} G. Weihs, T. Jennewein, C. Simon, H. Weinfurter, A. Zeilinger, ``Violation of Bell's inequality under strict Einstein
locality conditions'', \emph{Phys. Rev. Lett.} 81, 5039 (1998).
\bibitem{hensen} B. Hensen et al., ``Loophole-free Bell inequality violation using electron spins separated by 1.3 kilometres'',
\emph{Nature} 526, 682 (2015).
\bibitem{giustina} M. Giustina et al., ``Significant-loophole-free test of Bell's theorem with entangled photons'', \emph{Phys. Rev. Lett.}
115, 250401 (2015).
\bibitem{shalm} L.~K. Shalm et al., ``Strong loophole-free test of local realism'', \emph{Phys. Rev. Lett.} 115, 250402 (2015).
\bibitem{rosenfeld} W. Rosenfeld et al., ``Event-ready Bell test using entangled atoms simultaneously closing detection and locality
loopholes'', \emph{Phys. Rev. Lett.} 119, 010402 (2017).
\bibitem{storz} S. Storz et al., ``Loophole-free Bell inequality violation with superconducting circuits'', \emph{Nature} 617, 265 (2023).
\bibitem{handsteiner} J. Handsteiner et al., ``Cosmic Bell test: measurement settings from Milky Way stars'', \emph{Phys. Rev. Lett.}
118, 060401 (2017).
\bibitem{rauch} D. Rauch et al., ``Cosmic Bell test using random measurement settings from high-redshift quasars'', \emph{Phys. Rev.
Lett.} 121, 080403 (2018).
\bibitem{bigbell} The BIG Bell Test Collaboration, ``Challenging local realism with human choices'', \emph{Nature} 557, 212 (2018).
\bibitem{pr} S. Popescu, D. Rohrlich, ``Quantum nonlocality as an axiom'', \emph{Found. Phys.} 24, 379 (1994).
\bibitem{ic} M. Paw{\l}owski, T. Paterek, D. Kaszlikowski, V. Scarani, A. Winter, M. \.{Z}ukowski, ``Information causality as a physical
principle'', \emph{Nature} 461, 1101 (2009).
\bibitem{hall} M.~J.~W. Hall, ``Local deterministic model of singlet state correlations based on relaxing measurement independence'',
\emph{Phys. Rev. Lett.} 105, 250404 (2010).
\bibitem{hp} S. Hossenfelder, T. Palmer, ``Rethinking superdeterminism'', \emph{Front. Phys.} 8, 139 (2020).
\bibitem{brunner} N. Brunner, D. Cavalcanti, S. Pironio, V. Scarani, S. Wehner, ``Bell nonlocality'', \emph{Rev. Mod. Phys.} 86, 419 (2014).
\bibitem{pironio} S. Pironio et al., ``Random numbers certified by Bell's theorem'', \emph{Nature} 464, 1021 (2010).
\bibitem{barnum} H. Barnum, C.~M. Caves, C.~A. Fuchs, R. Jozsa, B. Schumacher, ``Noncommuting mixed states cannot be broadcast'',
\emph{Phys. Rev. Lett.} 76, 2818 (1996).
\bibitem{rt} S. Ryu, T. Takayanagi, ``Holographic derivation of entanglement entropy from AdS/CFT'', \emph{Phys. Rev. Lett.} 96, 181602
(2006).
\bibitem{adh} A. Almheiri, X. Dong, D. Harlow, ``Bulk locality and quantum error correction in AdS/CFT'', \emph{JHEP} 04 (2015) 163.
\bibitem{happy} F. Pastawski, B. Yoshida, D. Harlow, J. Preskill, ``Holographic quantum error-correcting codes: toy models for the
bulk/boundary correspondence'', \emph{JHEP} 06 (2015) 149.
\bibitem{bose} S. Bose et al., ``Spin entanglement witness for quantum gravity'', \emph{Phys. Rev. Lett.} 119, 240401 (2017).
\bibitem{mv} C. Marletto, V. Vedral, ``Gravitationally induced entanglement between two massive particles is sufficient evidence of
quantum effects in gravity'', \emph{Phys. Rev. Lett.} 119, 240402 (2017).
\bibitem{ggs} T.~D. Galley, F. Giacomini, J.~H. Selby, ``A no-go theorem on the nature of the gravitational field beyond quantum
theory'', \emph{Quantum} 6, 779 (2022).
\bibitem{azizhowl} J. Aziz, R. Howl, ``Classical theories of gravity produce entanglement'', \emph{Nature} 646, 813 (2025).
\bibitem{movw} C. Marletto, J. Oppenheim, V. Vedral, E. Wilson, ``Classical gravity cannot mediate entanglement'', arXiv:2511.07348
(2025).
\bibitem{diosi} L. Di\'osi, ``No, classical gravity does not entangle quantized matter fields'', arXiv:2511.00852 (2025).
\bibitem{shl} M.~D. Schneider, N. Huggett, N. Linnemann, ``A demonstration that classical gravity does not produce entanglement'',
arXiv:2511.19242 (2025).
\bibitem{oppenheim} J. Oppenheim, ``A postquantum theory of classical gravity?'', \emph{Phys. Rev. X} 13, 041040 (2023).
\bibitem{osswd} J. Oppenheim, C. Sparaciari, B. \v{S}oda, Z. Weller-Davies, ``Gravitationally induced decoherence vs space-time diffusion:
testing the quantum nature of gravity'', \emph{Nat. Commun.} 14, 7910 (2023).
\bibitem{ktm} D. Kafri, J.~M. Taylor, G.~J. Milburn, ``A classical channel model for gravitational decoherence'', \emph{New J. Phys.} 16,
065020 (2014).
\bibitem{westphal} T. Westphal, H. Hepach, J. Pfaff, M. Aspelmeyer, ``Measurement of gravitational coupling between millimetre-sized
masses'', \emph{Nature} 591, 225 (2021).
\bibitem{be} F. Englert, R. Brout, ``Broken symmetry and the mass of gauge vector mesons'', \emph{Phys. Rev. Lett.} 13, 321 (1964).
\bibitem{higgs} P.~W. Higgs, ``Broken symmetries and the masses of gauge bosons'', \emph{Phys. Rev. Lett.} 13, 508 (1964).
\bibitem{anderson} P.~W. Anderson, ``Plasmons, gauge invariance, and mass'', \emph{Phys. Rev.} 130, 439 (1963).
\bibitem{weinberg} S. Weinberg, ``A model of leptons'', \emph{Phys. Rev. Lett.} 19, 1264 (1967).
\bibitem{atlas} ATLAS Collaboration, ``Observation of a new particle in the search for the Standard Model Higgs boson with the ATLAS
detector at the LHC'', \emph{Phys. Lett. B} 716, 1 (2012).
\bibitem{cms} CMS Collaboration, ``Observation of a new boson at a mass of 125 GeV with the CMS experiment at the LHC'', \emph{Phys.
Lett. B} 716, 30 (2012).
\bibitem{pdgw} Particle Data Group, W mass world average, Review of Particle Physics 2024, \url{https://pdglive.lbl.gov/DataBlock.action?node=S043M}.
\bibitem{pdgn} Particle Data Group, neutron mean life, \url{https://pdglive.lbl.gov/view/S017T}.
\bibitem{pdgpi} Particle Data Group, charged pion mean life, \url{https://pdglive.lbl.gov/DataBlock.action?node=S008T}.
\bibitem{fewell} M.~P. Fewell, ``The atomic nuclide with the highest mean binding energy'', \emph{Am. J. Phys.} 63, 653 (1995).
\bibitem{superk} Super-Kamiokande Collaboration, ``Search for proton decay via $p\to e^+\pi^0$ and $p\to\mu^+\pi^0$ with an enlarged
fiducial volume in Super-Kamiokande I--IV'', \emph{Phys. Rev. D} 102, 112011 (2020).
\bibitem{borexino} Borexino Collaboration, ``Test of electric charge conservation with Borexino'', \emph{Phys. Rev. Lett.} 115, 231802
(2015).
\bibitem{planckxxiv} Planck Collaboration, ``Planck intermediate results. XXIV. Constraints on variations in fundamental constants'',
\emph{Astron. Astrophys.} 580, A22 (2015).
\bibitem{espresso} M.~T. Murphy et al., ``Fundamental physics with ESPRESSO: precise limit on variations in the fine-structure constant
towards the bright quasar HE 0515$-$4414'', \emph{Astron. Astrophys.} 658, A123 (2022).
\bibitem{bbn} I.~M. Dreyer, C.~J.~A.~P. Martins, ``Probing unification scenarios with big bang nucleosynthesis'', \emph{Phys. Rev. D} 113 (2026), arXiv:2604.04870.
\bibitem{oscillons} M.~A. Amin, R. Easther, H. Finkel, R. Flauger, M.~P. Hertzberg, ``Oscillons after inflation'', \emph{Phys. Rev.
Lett.} 108, 241302 (2012).
\bibitem{chou} C.~W. Chou, D.~B. Hume, T. Rosenband, D.~J. Wineland, ``Optical clocks and relativity'', \emph{Science} 329, 1630 (2010).
\bibitem{pr60} R.~V. Pound, G.~A. Rebka, Jr., ``Apparent weight of photons'', \emph{Phys. Rev. Lett.} 4, 337 (1960).
\bibitem{cassini} B. Bertotti, L. Iess, P. Tortora, ``A test of general relativity using radio links with the Cassini spacecraft'',
\emph{Nature} 425, 374 (2003).
\bibitem{israel} W. Israel, ``Event horizons in static vacuum space-times'', \emph{Phys. Rev.} 164, 1776 (1967).
\bibitem{eht} The Event Horizon Telescope Collaboration, ``First M87 Event Horizon Telescope results. I. The shadow of the supermassive
black hole'', \emph{Astrophys. J. Lett.} 875, L1 (2019).
\bibitem{hawking} S.~W. Hawking, ``Particle creation by black holes'', \emph{Commun. Math. Phys.} 43, 199 (1975).
\bibitem{page} D.~N. Page, ``Information in black hole radiation'', \emph{Phys. Rev. Lett.} 71, 3743 (1993).
\bibitem{lange} R. Lange et al., ``Improved limits for violations of local position invariance from atomic clock comparisons'',
\emph{Phys. Rev. Lett.} 126, 011102 (2021).
\end{thebibliography}

\end{document}
